% Privacy-first section 1. Budget 1.0 page, no figure.
\section{Introduction}
\label{sec:intro}

Genomic language models are released pretrained and fine-tuned locally, increasingly on cohort data
that the releasing laboratory never sees. A group with a few thousand participants' sequences adapts
a published checkpoint, reports a downstream metric, and shares the weights. We ask what that
sharing does to the people in the cohort.

A recent study established that published genomic models memorise their fine-tuning sequences to
very different degrees, without identifying what causes the difference
\citep{nemecek2026quantifying}. The natural follow-up is to find the cause, and we do; but the more
consequential finding turned out to be that the safety check the field applies to answer this
question is the wrong one, in a specific and repeatable way.

Our contributions, in the order the paper makes them:

\begin{enumerate}\itemsep2pt
\item \textbf{Every released masked checkpoint we tested returns planted records verbatim at a
held-out loss a practitioner would accept} (\S\ref{sec:leak}). The cheapest leak costs seven
thousandths of a bit and returns 30\% of records. On real chromosome-22 haplotypes --- each carried
by at most five of 2{,}548 individuals --- a record present as a single copy comes back letter for
letter roughly one time in ten.

\item \textbf{A model identifies an individual at checkpoints where it reproduces nothing at all}
(\S\ref{sec:ident}). Across three seeds and six checkpoints, single-copy verbatim recovery is exactly
zero in every cell, while identification of the same records rises to 0.92--0.96 against a chance
level of $0.0099$, with none of 350 never-trained controls ever identified. Reproduction and
identification have to be reported together; reporting either alone describes a different model than
the one that exists.

\item \textbf{What decides whether a leak is reachable is how many token slots the record occupies,
not how much the model knows} (\S\ref{sec:slots}). Across six cut sizes at matched held-out loss the
extraction rate is monotone in the slot count without exception, and the loss it costs to reach full
extraction falls in the same direction. Models at both ends answer any single hidden nucleotide
identically well; only reachability differs. The prediction this account makes about the released
DNABERT family --- no trend in $k$, because its overlapping tokens leave the slot count unchanged ---
was recorded before measuring and held.

\item \textbf{The number practitioners select on does not move} (\S\ref{sec:blind}). Downstream task
score changes by one to two per cent across the whole interval in which a model learns to identify
individuals, and on chromosome~22 its own maximum is a checkpoint that already leaks. Stopping where
held-out loss stops improving does work, with a margin of three hundredths of a bit.

\item \textbf{No single extraction test is correct across tokenizers} (\S\ref{sec:measure}). The same
checkpoint reads 0\% or 95\% depending only on which decoder was run, and the two protocols fail in
opposite directions. Had we used one throughout, we would have drawn the opposite conclusion about
tokenisation. This is why the measurement work belongs in the paper and not in an appendix.
\end{enumerate}

Two framings we deliberately avoid. We do not claim the character tokenizer is safe: HyenaDNA is a
released single-nucleotide model of the same size as our controlled backbone and it extracts 95\%,
so the tokenisation result is scoped to one architecture and one objective at a time. And we do not
report our decoder's failures as model safety --- where it reads zero we show, using a leave-one-out
bound, whether the content is nonetheless present.

All code, the measured output behind every figure, and the working logs are public
(\CODEURL).
